% Farben für jede Schicht (\providecommand: auch im Skript definiert)
\providecommand{\applicationLayerColor}{Gold}
\providecommand{\transportLayerColor}{Indigo}
\providecommand{\networkLayerColor}{RoyalBlue}
\providecommand{\dataLinkLayerColor}{Gray}

% command to make sure colorboxes don't create whitespace between code lines
\fboxsep=.6pt

% default settings for listings
\lstset{style=python}

\title[Netzwerke]{Netzwerke}

\subtitle{Schichtenmodell (Skript, Abschnitt 1.2)}

\author[Blum, Cyril] % (optional)
{Cyril Blum}

\institute[KLW] % (optional)
{
    Fachschaft Informatik\\
    Kantonsschule im Lee
}

\date[\today] % (optional)

\logo{\includegraphics[width=2cm]{Logos/FreeFlower.pdf}}

\titlegraphic{\vspace{-.7cm}\includegraphics[width=.85\textwidth]{"Grundlagen_Info/06_Datenbanken/Figures/title-emojis"}}

\frame{\titlepage}

% Eingabefelder für die Spielkarten (\providecommand: Foliensätze können gemeinsam eingebunden werden)
\providecommand{\fldNum}{\refstepcounter{FormCounter}\raisebox{-0.28\normalbaselineskip}{\TextField[name=\theFormCounter, width=2cm,print,borderwidth=0, backgroundcolor={0.9 0.9 0.9}]{}}}

\begin{frame}{Netzwerk-Spiel (ca. 5')}
    Sie sind ein Computer im Internet und möchten gerne eine Nachricht an einen anderen Computer senden:
    \begin{center}
        \begin{tcolorbox}[colframe=\networkLayerColor, colback=white, boxrule=1pt, arc=4pt, left=2pt, right=2pt, top=2pt, bottom=2pt]
            \begin{table}[H]
                \tiny
                \centering
                \begin{tabular}{ll}
                    \coin{3} Von:       & \texttt{192.168.0.\fldNum}                                                           \\\midrule
                    \coin{3} An:        & \texttt{192.168.0.\fldNum}                                                           \\\midrule
                    \coin{2} Sequenz:   & \texttt{\TextField[width=5cm, print,borderwidth=0, backgroundcolor={0.9 0.9 0.9}]{}} \\\midrule
                    \coin{1} Nachricht: & \TextField[width=5cm,print,borderwidth=0, backgroundcolor={0.9 0.9 0.9}]{}           \\
                \end{tabular}
            \end{table}
        \end{tcolorbox}
    \end{center}
    \textbf{Anleitung}\footnotesize
    \begin{itemize}[<+->]
        \item (Bereits gemacht): Notieren Sie sich eine einzigartige \acrshort{IP}-Adresse für Ihren Computer, startend mit \texttt{192.168.0.xyz} (z.B. \texttt{192.168.0.3}), wobei  $\texttt{xyz} \in [0,255]$
        \item[\coin{1}] (Vorderseite) Schreiben Sie eine Nachricht (max. 10 Zeichen pro Karte)
        \item[\coin{2}] (Vorderseite) Nummerieren Sie die Karten in der Reihenfolge, in der sie gesendet werden sollen (Sequenz)
        \item[\coin{3}] (Rückseite) Fügen Sie die \acrshort{IP}-Adressen der Absender- und Empfänger-Computer hinzu.
        \item[\faTruck] Übergeben Sie die Karten dem \enquote{Router} (= \enquote{Postbote})
    \end{itemize}
\end{frame}

\begin{frame}[fragile]{Was sind Netzwerke?}
    \begin{itemize}[<+->]
        \item Graphen:
              \begin{figure}[H]
                  \centering
                  \includegraphics[keepaspectratio, width=0.35 \paperwidth]{Grundlagen_Info/02_Graphen/Figures/4-graph-zusammenhaengend.png}
              \end{figure}

        \item Im Alltag:
              \begin{itemize}
                  \item Kommunikations-Netzwerke: Die Post, DHL, Planzer, ...
                  \item Soziale Netzwerke: TikTok, YouTube, Instagram, ...
                  \item Transport-Netzwerke: SBB, Flixbus, ...
                  \item Shopping-Netzwerke: Aliexpress, Uber, ...
                  \item Computer-Netzwerke: \textbf{Das Internet}
              \end{itemize}
    \end{itemize}
\end{frame}

\begin{frame}[fragile]{Schichtenmodell}{Analogie}
    \begin{figure}[H]
        \centering
        \begin{tikzpicture}[
                auto,
                block/.style={
                        minimum width=3cm,
                        node distance=1.5cm,
                        rectangle,
                        fill=\networkLayerColor!50,
                        align=center,
                        rounded corners
                    }
            ]


            % Define the nodes with FontAwesome icons
            \node[block] (chefs) {\faFemale\ Chef};
            \node[block, below=of chefs] (secretariat) {\faPhone\ Sekretariat};
            \node[block, below=of secretariat] (postman) {\faMale\ Post};


            % Define the nodes with FontAwesome icons

            \node[block, right=of postman] (postman2) {\faMale\ Post};
            \node[block, above=of postman2] (secretariat2) {\faPhone\ Sekretariat};
            \node[block, above=of secretariat2] (chefs2) {\faMale\ Chef};


            % Draw the dashed arrows
            \draw[<->, dashed, \networkLayerColor!50, thick] (chefs) -- (chefs2);
            \draw[<->, dashed, \networkLayerColor!50, thick] (secretariat) -- (secretariat2);
            \draw[<->, dashed, \networkLayerColor!50, thick] (postman) -- (postman2);

            \draw[->, line width=1mm,ForestGreen!50] (chefs) -- node[left]{\faComments} (secretariat) ;
            \draw[->, line width=1mm,ForestGreen!50] (secretariat) -- node[left]{\faEnvelope} (postman) ;
            \draw[->, line width=1mm,ForestGreen!50, bend right] (postman.south) to node[midway,below]{\faBicycle\ \faMotorcycle\ \faCar\ \faShip\ \faPlane} (postman2.south) ;
            \draw[->, line width=1mm,ForestGreen!50] (postman2) -- node[right]{\faEnvelope} (secretariat2) ;
            \draw[->, line width=1mm,ForestGreen!50] (secretariat2) -- node[right]{\faComments} (chefs2) ;

            % fit rectangles around regions
            \node[fit=(chefs) (postman), inner sep=8pt, red, dotted, draw, line width=1pt,rounded corners](region1) {};
            \node[fit=(chefs2) (postman2), inner sep=8pt, red, dotted, draw, line width=1pt,rounded corners](region2) {};
            % annotate both rectangles
            \node[above=5pt of region1.north, anchor=center, align=center]{\textcolor{red}{Firma 1, Europa}};
            \node[above=5pt of region2.north, anchor=center, align=center]{\textcolor{red}{Firma 2, Asien}};

        \end{tikzpicture}
    \end{figure}
\end{frame}


\begin{frame}[fragile]{Schichtenmodell: Übersicht}
    \Wider[3em]{\begin{center}
            \begin{table}[H]
                \centering
                \footnotesize
                \adjustbox{max width=\textwidth, max totalheight=.78\textheight}{
                    \begin{tabular}{cm{2.6cm}m{2cm}m{4.6cm}m{2.5cm}m{2cm}}
		\textbf{\#}              & \textbf{Schicht (\acrshort{OSI})}                     & \textbf{Schicht \newline (\acrshort{TCP} / \acrshort{IP})}  & \textbf{Aufgabe}                                                                                                & \textbf{Informations\-form}               & \textbf{Begriffe / \newline \underline{Protokolle}}   \\ \midrule
		\rowcolor{\applicationLayerColor!10}7    & Anwendung \newline \textit{Application}    & Anwendung\newline\textit{Application} & Hilft bei der Identifizierung des Clients und synchronisiert die Kommunikation.                                 & Nachricht                                 & \underline{\acrshort{HTTP}}, \underline{\acrshort{HTTPS}}, \underline{\acrshort{DNS}}, \underline{\acrshort{SMTP}}, \underline{\acrshort{FTP}}...        \\\midrule
		\rowcolor{\applicationLayerColor!20}6    & Darstellung \newline \textit{Presentation} & Anwendung\newline\textit{Application} & Daten von der Anwendungsschicht werden extrahiert und für die Übertragung in das erforderliche Format gebracht. & Nachricht                                 & \underline{\acrshort{TLS}}, \underline{\acrshort{SSL}}, \acrshort{ASCII}, \acrshort{UTF-8}, ...      \\\midrule
		\rowcolor{yellow!30}5    & Sitzung \newline \textit{Session}          & Anwendung\newline\textit{Application} & Stellt Verbindung her, hält diese aufrecht, gewährleistet Authentifizierung und sichert die Sicherheit.         & Nachricht (oder verschlüsselte Nachricht) & \underline{\acrshort{TLS}} Handshake           \\\midrule
		\rowcolor{\transportLayerColor!20}4    & Transport\newline \textit{Transport}       & Transport\newline\textit{Transport}   & Nimmt Dienst von der Vermittlungsschicht und stellt ihn der Anwendungsschicht zur Verfügung.                    & Segment                                   & \underline{\acrshort{TCP}}, \underline{\acrshort{UDP}}, Ports... \\\midrule
		\rowcolor{\networkLayerColor!20}3 & Vermittlung \newline \textit{Network}      & Internet\newline\textit{Internet}     & Übertragung von Daten von einem Host zu einem anderen, der sich in unterschiedlichen Netzwerken befindet.       & Paket                                     & \underline{\acrshort{IP}v4}, \underline{\acrshort{IP}v6}, \underline{\acrshort{ICMP}}, Router, Gateway   \\\midrule
		\rowcolor{\dataLinkLayerColor!30}2      & Sicherung \newline \textit{Data Link}      & Netzzugangs\-schicht\newline\textit{Link /\newline Network\newline Access} & Übermittlung von Nachrichten von Knoten zu Knoten.                                                              & Frame                                     & \underline{\acrshort{ARP}}, \acrshort{MAC}, Switch, Bridge, Ethernet, \acrshort{WiFi}, ...      \\\midrule
		\rowcolor{\dataLinkLayerColor!15}1      & Bitübertragung \newline \textit{Physical}  & Netzzugangs\-schicht\newline\textit{Link /\newline Network\newline Access} & Herstellung physischer Verbindungen zwischen Geräten.                                                           & Bits                                      & Hub, Repeater, Modem, Kabel...        \\
	\end{tabular}}
            \end{table}
        \end{center}}
\end{frame}


\setbeamercolor{frametitle}{bg=white}

\begin{frame}[fragile]{Schichten \& Protokolle}
    \begin{figure}[H]
        \centering
        \adjustbox{max width=.8\textwidth, max height=.35\textheight}{
            \begin{tikzpicture}[
        envelope/.style={
                draw,
                thick,
                rounded corners=3pt,
                font=\small,
                align=center
            },
        layerlabel/.style={
                font=\scriptsize,
                anchor=north west
            },
        brace/.style={decorate, decoration={brace}, red, thick},
        rbrace/.style={decorate, decoration={brace, mirror}, red, thick}
    ]

    % ---- Sizes per layer ----
    \def\increaseW{1cm}
    \def\increaseH{1cm}
    \def\moveaway{8pt}
    \def\wA{5cm}   \def\hA{.8cm}   % Anwendung
    \def\wT{\wA+\increaseW}   \def\hT{\hA+\increaseH}   % Transport
    \def\wV{\wT+\increaseW}   \def\hV{\hT+\increaseH}   % Vermittlung
    \def\wN{\wV+\increaseW}   \def\hN{\hV+\increaseH}   % Netzzugangsschicht

    % ---- Netzzugangsschicht ----
    \node[envelope, fill=\dataLinkLayerColor!15, minimum width=\wN, minimum height=\hN] (net) {};
    \node[layerlabel] at (net.north west) {Netzzugangsschicht (\texttt{\acrshort{MAC}} zu \texttt{\acrshort{MAC}})};

    % ---- Vermittlung ----
    \node[envelope, fill=\networkLayerColor!20, minimum width=\wV, minimum height=\hV] (ver)
    at ($(net.center)+(0,-1/2*\increaseH)$) {};
    \node[layerlabel] at (ver.north west) {Vermittlung (\texttt{\acrshort{IP}} zu \texttt{\acrshort{IP}})};

    % ---- Transport ----
    \node[envelope, fill=\transportLayerColor!20, minimum width=\wT, minimum height=\hT] (tr)
    at ($(net.center)+(0,-2/2*\increaseH)$) {};
    \node[layerlabel] at (tr.north west) {Transport (Port zu Port)};

    % ---- Anwendung ----
    \node[envelope, fill=\applicationLayerColor!20, minimum width=\wA, minimum height=\hA] (app)
    at ($(net.center)+(0,-3/2*\increaseH)$) {};
    \node[layerlabel] at (app.north west) {Anwendung (Client zu Server)};

    % ---- Braces for Header and Payload ----
    % Payload (Nutzteil): only the application data (inner payload)
    \draw[brace] ($(app.north east)+(3/2*\increaseH+4pt,0)$) -- ($(app.south east)+(3/2*\increaseH+4pt,0)$)
    node[pos=0.5, align=left, font=\small, right=4pt] {\textbf{Nutzteil}\\(Payload)};

    % dotted lines for payload
    \draw[red, dotted, thick] (app.north east) -- ($(app.north east)+(3/2*\increaseH+\moveaway,0)$);
    \draw[red, dotted, thick] (app.south east) -- ($(app.south east)+(3/2*\increaseH+\moveaway,0)$);

    % Header (Kopfteil): all headers from all layers
    \draw[rbrace] ($(net.north west)-(\moveaway,0)$) -- ($(app.north west)+(-3/2*\increaseH-\moveaway,0)$)
    node[pos=0.5, align=right, font=\small, left=4pt] {\textbf{Kopfteil}\\(Header)};

    % dotted lines for header
    \draw[red, dotted, thick] (tr.south west) -- ($(app.north west)+(-3/2*\increaseH-\moveaway,0)$);
    \draw[red, dotted, thick] (net.north west) -- ($(net.north west)-(\moveaway,0)$);
\end{tikzpicture}

        }
    \end{figure}
    \small\pause
    \only<2>{
        \begin{table}[H]
            \centering
            \begin{tabular}{p{2cm}l}
                \rowcolor{\networkLayerColor!20} Von:                                         & \texttt{192.168.0.3}  \\\midrule
                \rowcolor{\networkLayerColor!20} An:                                          & \texttt{192.168.0.11} \\\midrule
                \rowcolor{\transportLayerColor!20}Sequenz:                                    & 1                     \\\midrule
                \rowcolor{\applicationLayerColor!20}Nachricht\newline\tiny (max. 10 Zeichen): & \texttt{Hi Bob, wi}   \\
            \end{tabular}
            \begin{tabular}{p{2cm}l}
                \rowcolor{\networkLayerColor!20} Von:                                         & \texttt{192.168.0.3}  \\\midrule
                \rowcolor{\networkLayerColor!20} An:                                          & \texttt{192.168.0.11} \\\midrule
                \rowcolor{\transportLayerColor!20}Sequenz:                                    & 2                     \\\midrule
                \rowcolor{\applicationLayerColor!20}Nachricht\newline\tiny (max. 10 Zeichen): & \texttt{e geht's?}    \\
            \end{tabular}
        \end{table}
    }
    \pause
    % onlyenv: the columns take no space on slide 2 (where the two tables are shown)
    \begin{onlyenv}<3->
    \begin{columns}[t]
        \begin{column}{.5\textwidth}
            \textbf{Schichten}
            \begin{itemize}[<+->]
                \item Z.B. Anwendungs-, Transport-, Vermittlungs-, Netzzugangsschicht
                \item Erlauben es, spezialisierte Aufgaben in einzelne Schichten aufzuteilen: Z.B. Übermittlung von Signalen, Sprache-Regeln etc.
            \end{itemize}
        \end{column}
        \pause
        \begin{column}{.5\textwidth}
            \textbf{Protokolle}
            \begin{itemize}[<+->]
                \item Definiert Ablauf einer Kommunikation, Aufbau der Datenpakete
                \item Gibt es auf jeder Ebene des Schichtenmodells
                \item Analogie: Begrüssungsformen, \textit{small talk}, etc.
            \end{itemize}
        \end{column}
    \end{columns}
    \end{onlyenv}
\end{frame}

\begin{frame}[fragile]{Schichtenmodell}{Vorteile}
    \begin{itemize}[<+->]
        \item \textbf{Modularität}: Aufbau auf bestehenden Schichten
        \item \textbf{Austauschbarkeit} / Flexibilität
              \begin{itemize}
                  \item Übertragung über Kabel, \acrshort{WLAN} etc. kann ausgetauscht werden, ohne dass darüberliegende Schichten tangiert werden
                  \item Analogie: Anstelle von Lastwagen kann Motorrad verwendet werden
              \end{itemize}
    \end{itemize}
\end{frame}


\begin{frame}{Netzwerk-Spiel}{Diskussion}
    \begin{itemize}[<+->]
        \item Was sind die Aufgaben der einzelnen Schichten?
              \begin{itemize}
                  \item Anwendungsschicht
                  \item Transportschicht
                  \item Internetschicht
                  \item Netzzugangsschicht
              \end{itemize}
        \item Sind Pakete immer angekommen?
        \item Wo waren die Flaschenhälse bei der Übermittlung der Nachrichten?
        \item Was ist der Unterschied zwischen einem Paket und einer Nachricht?
        \item \aufgabebox{1.1 (Netzwerkspiel): Schichten, Protokolle, Geräte}
    \end{itemize}
\end{frame}



\begin{frame}[fragile]{Schichtenmodell}{Video}
    \begin{figure}[H]
        \href{https://www.youtube.com/watch?v=PBWhzz_Gn10}{YouTube-Video (13')}
    \end{figure}

    \begin{itemize}
        \item \aufgabebox{1.43 (Video: Schichten, Geräte, Protokolle notieren)}
        \item \aufgabebox{1.44 (Austausch in 3er- bis 4er-Gruppen, ca. 5')}
    \end{itemize}
\end{frame}

